\documentclass[11pt]{article}
\usepackage[margin=1in]{geometry}
\usepackage{amsmath, amssymb, amsthm}
\usepackage[most]{tcolorbox}

\definecolor{definitionbody}{HTML}{E5FFE5}
\definecolor{definitiontitle}{HTML}{C5FFC5}
\definecolor{theorembody}{HTML}{E3F4FA}
\definecolor{theoremtitle}{HTML}{C5E7F2}

\tcbset{
    notesbox/.style={
        enhanced,
        breakable,
        frame hidden,
        boxrule=0pt,
        arc=2pt,
        coltitle=black,
        fonttitle=\bfseries,
        left=10pt,
        right=10pt,
        top=7pt,
        bottom=7pt,
        toptitle=5pt,
        bottomtitle=5pt,
        before skip=12pt,
        after skip=12pt
    }
}

\newtcolorbox{definition}[1][Definition]{
    notesbox,
    title={#1},
    colback=definitionbody,
    colbacktitle=definitiontitle
}

\newtcolorbox{theorem}{
    notesbox,
    title=Theorem,
    colback=theorembody,
    colbacktitle=theoremtitle
}

\newtcolorbox{corollary}{
    notesbox,
    title=Corollary,
    colback=theorembody,
    colbacktitle=theoremtitle
}

\title{Section 19 Notes}
\author{Jonathan Nguyen}
\date{}

\begin{document}
\maketitle

\begin{definition}[Divisors of Zero]
If $a,b \in R$ are nonzero and $ab=0$, then
$a$ and $b$ are called divisors of zero.
\end{definition}

\begin{theorem}
In the ring $\mathbb{Z}_n$, the divisors of zero are
precisely those nonzero elements that are not
relatively prime to $n$.
\end{theorem}

\noindent
Note: Find the nonzero elements $a \in \mathbb{Z}_n$
such that $\gcd(a,n) \neq 1$.

\begin{corollary}
If $p$ is prime, then $\mathbb{Z}_p$ has no
divisors of zero.
\end{corollary}

\begin{definition}[Integral Domain]
An integral domain $D$ is a commutative ring
with unity $1 \neq 0$ and containing no
divisors of zero.
\end{definition}

\begin{theorem}
Every field $F$ is an integral domain.
\end{theorem}

\begin{theorem}
Every finite integral domain is a field.
\end{theorem}

\begin{definition}[Characteristic of a Ring]
The characteristic of a ring $R$ is the least
positive integer $n$ such that $na=0$ for all
$a \in R$. If no such positive integer exists,
the characteristic of $R$ is $0$.
\end{definition}

\begin{theorem}
Let $R$ be a ring with unity $1$.
If $n \cdot 1 \neq 0$ for every
$n \in \mathbb{Z}^{+}$, then $R$ has
characteristic $0$.
\end{theorem}

\end{document}